% ------------------------------------------------------------
% Page 108
% ------------------------------------------------------------

\textbf{« ARGELIERS »} le mari de Mme. Argeliers, la casquette de travers, long « comme
un jour sans pain » m’avait pris en amitié. C’était un piégeur professionnel de fouines
et de martres dont les peaux se négociaient à Millau, au prix fort de 6 à 800F. de
l’époque. Il me proposa de l’accompagner ne voyant pas en moi un « concurrent
possible ! ». Je le suivis dans des endroits impossibles dans le silence de ses pas
feutrés ; sans parole, il me montrait du doigt l’emplacement où se trouvaient les
pièges, mais je ne distinguais rien ; lui, à d’infinis détails, il voyait et savait. Ce jour-
là nous revînmes bredouilles mais il me confia que les mois d’hiver il gagnait pour
toute l’année.

Les beaux jours revinrent ; durant le mois de mai on installa l’électricité et je fus
surpris d’apprendre que la maison d’école ne figurait pas sur le projet. Je crus donc,
par devoir vis-à-vis de mon possible successeur, d’alerter encore une fois M. le
Maire... La réponse fut brève : « La Mairie de Meyrueis n’est en rien concernée. » Je
partis en claquant la porte.

\begin{center}
\textbf{Juin-juillet : Fin de l’année scolaire}
\end{center}

\textbf{Juin :}

Nous faisions de longues promenades dans les forêts domaniales, scrutant les fossés,
espérant découvrir morilles ou quelques cèpes, mais ce printemps-là, trop sec, n’était
pas favorable. Un jour, longeant un sentier remontant la rivière, nous fûmes étonnés
d’atteindre si vite le col de la Serreyrède et là, signalé à 2km, le village de l’Espérou
qui avait déjà une vocation touristique. Sur notre lancée nous arrivâmes au village.
C’était le jour de la fête votive ; nous y participâmes quelques heures avant de
reprendre le chemin du retour qui fut, par les pentes, très rapide, réalisant que
pédestrement nous étions plus proches de l’Espérou que de Meyrueis.

\textbf{Juillet :}

Notre départ était proche. Comment pouvions-nous remercier les gens de leur bon
accueil, de leur aide prodiguée gratuitement en bien des circonstances ?
L’idée nous vint de réunir au temple nos amis, c’est à dire tout le monde non pas pour
y entendre « la bonne parole », mais pour partager avec nous de modestes
« festivités ». Encore une fois nous déplaçames tables et bancs, on nous prêta tables et
chaises, vaisselle... car tout le monde voulait participer, le « Caïffa » nous livra, sur
commande, tout ce qui était nécessaire pour organiser un « buffet » honorable. Tout le
monde répondit à notre invitation : la soirée se prolongea fort tard et nous fûmes
heureux.
